\documentclass[12pt,a4paper,oneside]{article}
\usepackage[brazil]{babel}
\usepackage[utf8]{inputenc}
\usepackage{olymp}
\usepackage{color}
\usepackage[dvipsnames]{xcolor}
\usepackage{listings}

\setcounter{problem}{0}

\begin{document}

\begin{problem}{Conversor de horas}{convhoras}{2}

Faça um programa que leia um valor inteiro representando o tempo em segundos de um determinado evento e informe-o expresso no formato horas:minutos:segundos.

Para resolver o problema implemente obrigatoriamente uma função única que receba a duração total do evento em segundos e retorne por referência a quantidade de horas,
minutos e segundos em que esse valor total é decomposto, de acordo com o seguinte protótipo:

\texttt{\textcolor{YellowGreen}{\textbackslash\textbackslash Converte seg segundos em h horas m minutos s segundos}}\\
\texttt{\textcolor{Mulberry}{void} convhoras(\textcolor{Mulberry}{int} seg, \textcolor{Mulberry}{int \&}h, \textcolor{Mulberry}{int \&}m, \textcolor{Mulberry}{int \&}s)};




%\Notes
%
%Observações, se tiver.

\InputFile

A entrada consiste em um único valor inteiro $N$, representando uma quantidade total de segundos. Restrição: $0 \leq N \leq 86400$ (a quantidade de segundos em um dia).

\OutputFile

Seu programa deve gerar uma única linha contendo o valor convertido para horas, minutos e segundos, no formato ``\texttt{hh:mm:ss}'', ou seja, os valores devem ser separados por ``\texttt{:}'' e cada um deve ter dois dígitos (assim, valores de 0 a 9 devem ser impressos com um zero à esquerda).


% \Notes
% Nesta questão, seu programa deve usar a função \texttt{main} dada acima exatamente como está, não a modifique! Sua tarefa é apenas criar a função \texttt{eh\_primo}.

\Example

\begin{example}
\exmp{
556
}{
00:09:16
}%
\end{example}

\begin{example}
\exmp{
6666
}{
01:51:06
}%
\end{example}

%Comentários sobre o exemplo, se necessário.

\end{problem}

\end{document}
